\section{Methodology}
\label{sec:methodology}

\subsection{Task Graph Generation and Motifs}
To ensure structural diversity and avoid bias toward any single graph topology, we generate synthetic DAG task environments across three representative structural motifs (Figure~\ref{fig:topologies}):
\begin{enumerate}
    \item \textbf{Fork-Join Motif:} A wide parallel structure where an initial data generator fans out into two distinct processing pipelines (Branch A and Branch B) that subsequently merge into an intermediate join node, fan out into a secondary computation stage, and converge into a final join prior to terminal synthesis.
    \item \textbf{Spine-Join Motif:} An asymmetric structure featuring an iterative backbone (spine) running through the workflow, from which parallel auxiliary branches periodically detach, perform multi-step side-computations, and re-integrate into the main trunk via join operations.
    \item \textbf{Hierarchical Reduction Motif:} A balanced multi-tier binary tree reduction where pairs of parallel leaf computations recursively aggregate through successive join stages until reaching a single terminal summary.
\end{enumerate}

Every generated graph enforces the structural guarantee that at least two multi-parent dependency join nodes ($|\text{Parents}(v)| \ge 2$) exist, providing valid, non-trivial dependency integration boundaries for milestone scheduling across all horizon levels $H \in \{12, 24, 36\}$.

\subsection{Deterministic Verification Contracts}
Rather than relying on non-deterministic LLM-as-judge prompts that introduce stochastic evaluation noise, our verification layer implements a suite of formal deterministic contracts $\mathcal{C}$:
\begin{itemize}
    \item \textbf{Schema Conformance Contract:} Verifies that every intermediate dictionary artifact contains all mandatory structural fields (e.g., \texttt{records}, \texttt{schema}, \texttt{metadata}, \texttt{summary}) required by downstream consumers.
    \item \textbf{Type Integrity Contract:} Enforces strict runtime data types (e.g., verifying that dataset collections remain valid Python lists, record identifiers remain strings, and statistical metrics remain numerical floats).
    \item \textbf{Relational Reference Integrity Contract:} For multi-parent join operations, asserts that every record key in the joined artifact originates from the reference key set of at least one valid parent dependency:
    \begin{equation}
        \forall r \in x_{\text{join}}.\text{records}, \quad r.\text{id} \in \bigcup_{p \in \text{Parents}(v)} x_p.\text{keys}
    \end{equation}
    \item \textbf{Numeric Bound Contract:} Enforces domain-specific boundary constraints (e.g., normalized value intervals $[0.0, 100.0]$ and non-negative variances).
\end{itemize}

\subsection{Verification Scheduling Policies}
We implement and contrast four verification scheduling policies under equivalent execution conditions:

\paragraph{Arm A: Final-Only Verification ($\mathcal{S}_A$)}
Represents standard unverified execution. Verification is invoked exclusively at workflow completion:
\begin{equation}
    \mathcal{S}_A = \{H\}, \quad K_{\text{alloc}} = 1
\end{equation}

\paragraph{Arm B: Fixed-Periodic Verification ($\mathcal{S}_B$, Baseline)}
Allocates $K = 3$ verification slots at uniform step intervals:
\begin{equation}
    \mathcal{S}_B = \left\{ \left\lceil \frac{H}{3} \right\rceil, \left\lceil \frac{2H}{3} \right\rceil, H \right\}, \quad K_{\text{alloc}} = 3
\end{equation}
Check timing is determined solely by execution step count, completely independent of graph topology.

\paragraph{Arm C: Dependency-Milestone Verification ($\mathcal{S}_C$, Intervention)}
Allocates $K = 3$ verification slots immediately prior to dependency integration boundaries:
\begin{equation}
    \mathcal{S}_C = \{t_{\text{join1}}, t_{\text{join2}}, H\}, \quad K_{\text{alloc}} = 3
\end{equation}
where $t_{\text{join1}}$ is the execution step immediately preceding the execution of the first multi-parent join node (validating that all inputs to Join 1 have successfully completed), $t_{\text{join2}}$ is the step immediately preceding the second join node, and $H$ is the terminal step. This schedule directly tests the hypothesis that verifying intermediate state at dependency integration boundaries provides optimal error containment.

\paragraph{Arm D: Every-Step Verification ($\mathcal{S}_D$, Upper Bound)}
Invokes verification after every node execution:
\begin{equation}
    \mathcal{S}_D = \{1, 2, \dots, H\}, \quad K_{\text{alloc}} = H
\end{equation}
Serves as the empirical upper bound on early detection and error containment, establishing the maximum performance achievable with unlimited verification resources.

\subsection{Controlled Fault Injection Architecture}
Faults are injected dynamically during execution using an external \texttt{FaultInjector}. Each fault injection episode is parameterized by a formal \texttt{FaultDescriptor} specifying:
\begin{itemize}
    \item \textbf{Target Node ($v_{\text{fault}}$):} Chosen systematically within one of three horizon-relative position bins: \emph{Early} ($t \in [1, \lfloor H/3 \rfloor]$), \emph{Middle} ($t \in [\lfloor H/3 \rfloor + 1, \lfloor 2H/3 \rfloor]$), or \emph{Late} ($t \in [\lfloor 2H/3 \rfloor + 1, H]$).
    \item \textbf{Fault Type:} One of four controlled failure modes:
    \begin{enumerate}
        \item \texttt{missing\_field}: Drops a mandatory top-level key from the node output dictionary.
        \item \texttt{type\_corruption}: Mutates a record collection into an unhashable primitive type (e.g., replacing a list with an integer).
        \item \texttt{numeric\_corruption}: Multiplies numerical values by an extreme scaling factor, violating physical range boundaries.
        \item \texttt{broken\_reference}: Corrupts record identifiers, generating synthetic orphaned foreign keys.
    \end{enumerate}
\end{itemize}
Fault injection is strictly \textbf{one-shot}: the fault injector mutates state at $v_{\text{fault}}$ exactly once during the initial forward pass, isolating single-cause error propagation.

\subsection{Execution Engine and Recovery Policy}
Execution is governed by a unified \texttt{WorkflowExecutor} managing an immutable state snapshot store. Prior to each node execution, a lightweight state checkpoint is registered. 

When a scheduled verification check fails at step $t_{\text{detect}}$, the executor executes a pre-registered, uniform recovery policy:
\begin{enumerate}
    \item \textbf{State Rollback:} The executor discards all tainted execution state produced between $t_{\text{checkpoint}}$ and $t_{\text{detect}}$, restoring the most recent verified passing snapshot $x_{t_{\text{checkpoint}}}$.
    \item \textbf{Suffix Replay:} The executor re-executes the unexecuted workflow suffix from $t_{\text{checkpoint}}$ to $H$.
    \item \textbf{Fault Disarming:} During the replay pass, the fault injector is disarmed. This models a successful single retry or prompt refinement, allowing us to measure pure re-execution workload ($\text{ReplayNodes}$) without confounding recovery heuristics.
    \item \textbf{Check Invariance:} In accordance with pre-registered budget fairness protocols, no additional verification checks are invoked during the replay pass, ensuring that rollback does not inflate realized verifier budgets.
\end{enumerate}
